\begin{helper}[Replacement free-color full path]
With $\gamma,\varphi_\delta,T_\delta,X_\delta$, and $Z$ as above, there
are a color $\beta$, vertices $v\in Z$ and $z$, and a simple full
$(\gamma,\beta)$ incidence path $P_\delta$ under $\varphi_\delta$ from
$v$ to $z$ such that $\beta$ is missing at $v$ under both
$\varphi_\delta$ and $\varphi$, is free for $T_\delta$, differs from
both $\delta$ and $\gamma$, and $P_\delta$ has precisely the support of
the two-color component through $v$ and ends outside $X_\delta$.
\end{helper}
\begin{proof}
The replacement theorem \noderef{ScheideReplacementMaximum} makes
$(T_\delta,e_0,\varphi_\delta)$ globally maximum, and
\noderef{ScheideGlobalMaximumOrder} makes $X_\delta$ elementary and
closed.  Its exact vertex set is $Z\cup U_\delta$ by
\noderef{ScheideReplacementData}; in particular, the two ends of $e_0$
remain in the base $Y\subseteq Z$.

Apply \noderef{ScheideGlobalMaximumMissingCounts} to every vertex of
$Z$.  As in the original count, each vertex contributes at least one
missing color, the two ends of $e_0$ contribute one additional color
each, and elementariness makes the contributions at distinct vertices
disjoint.  Thus at least $p+2s$ distinct colors are missing on $Z$.
At most $p-2+s$ colors are used on $T_\delta$: at most $p-2$ come from
retained colored base edges, $s-1$ from retained blocks, and the two
appended edges have the single common color $\delta$ by
\noderef{ScheideReplacementPairFixed} and
\noderef{ScheideReplacementData}.

Every used color has a missing witness in $Z$.  A retained base or block
color has its old base witness from
\noderef{ScheideStructuralConfiguration}.  No retained displayed edge
has color $\delta$, because \noderef{ScheideGlobalMaximumUnused} says
that the defective color is unused on the original tree.  The switch
changes at an old
tree vertex only the two colors $\gamma,\delta$ at $v^0_\delta$, by
\noderef{ScheideDefectiveSwitchEffect}; retained base colors avoid all
block colors, retained block colors differ from $\gamma$ by injectivity,
and so those old witnesses survive.  The new used color $\delta$ is
missing at $v^0_\delta\in Y\subseteq Z$ by the exact missing-set formula
in the same switch-effect theorem.

Since $p-2+s<p+2s$, choose a color $\beta$ missing in $Z$ but unused on
$T_\delta$.  It is free by \noderef{TreeFreeColors}.  It differs from
$\delta$, which is used on the appended pair.  It differs from $\gamma$
as well: \noderef{ScheideSymmetricReplacementBoundaryConfiguration}
makes $\gamma$ defective, and closedness then makes $\gamma$ present at
every vertex of $X_\delta$.  The switch-effect formula, together with
$\beta\ne\gamma,\delta$, shows that the same vertex $v\in Z$ also misses
$\beta$ under the original coloring.

Finally, the globally maximum replacement triple, the defective-color
clause just proved, and this free missing color satisfy all hypotheses of
\noderef{ScheideDefectiveChainPathExists}.  That theorem supplies the
outside endpoint $z$ and the required full simple path, with the exact
support assertion expressed by \noderef{ScheideDefectiveChainPath}.
\end{proof}
